\documentclass[11pt]{article}
\usepackage{cornell-notes}
\usepackage{mathtools}
\usepackage{needspace}
\usepackage{ragged2e}
\tcbuselibrary{most}
\geometry{letterpaper,margin=0.72in,headheight=15pt}
\setCornellCollection{Cloud Architecture Cornell Notes}
\setCornellUnitType{Chapter}
\setCornellUnitNumber{2}
\setCornellUnitTitle{Case Study Background}

\definecolor{Navy}{HTML}{17324D}
\definecolor{Teal}{HTML}{147D78}
\definecolor{Sky}{HTML}{E8F2F7}
\definecolor{CueGray}{HTML}{EEF2F4}
\definecolor{Warm}{HTML}{FFF3D6}
\definecolor{Risk}{HTML}{953B35}
\definecolor{Green}{HTML}{276749}
\definecolor{RuleGray}{HTML}{AAB8C2}
\definecolor{Text}{HTML}{25313A}

\hypersetup{colorlinks=true,linkcolor=Navy,urlcolor=Teal,citecolor=Navy,pdfborder={0 0 0}}
\color{Text}
\setlength{\parindent}{0pt}
\setlength{\parskip}{4pt}
\setlength{\emergencystretch}{3em}
\setlength{\LTpre}{3pt}
\setlength{\LTpost}{5pt}
\renewcommand{\arraystretch}{1.16}
\setlist{nosep,leftmargin=1.45em}
\hfuzz=0.2pt
\hbadness=2500

\titleformat{\section}{\Large\bfseries\color{Navy}\RaggedRight}{}{0pt}{}[\color{Teal}\titlerule]
\titleformat{\subsection}{\large\bfseries\color{Teal}\RaggedRight}{}{0pt}{}
\titlespacing*{\section}{0pt}{1.35ex plus .4ex}{.7ex}
\titlespacing*{\subsection}{0pt}{1.1ex plus .3ex}{.45ex}

\pagestyle{fancy}
\fancyhf{}
\fancyhead[L]{\small\color{Navy}\textbf{Cloud Architecture Cornell Notes}}
\fancyhead[R]{\small\color{Teal}\ChapterMark}
\fancyfoot[L]{\footnotesize\color{Navy}Concepts $\bullet$ Mechanisms $\bullet$ Evidence $\bullet$ Decisions}
\fancyfoot[R]{\footnotesize\color{Navy}\thepage\ of \pageref*{LastPage}}
\renewcommand{\headrulewidth}{0.5pt}
\renewcommand{\footrulewidth}{0.35pt}

\newtcolorbox{orientationbox}{enhanced,breakable,colback=Sky,colframe=Teal,boxrule=.75pt,arc=2mm,title={Chapter Orientation},fonttitle=\bfseries,left=2mm,right=2mm,top=1.5mm,bottom=1.5mm}
\newtcolorbox{sectionsummarybox}{enhanced,breakable,colback=CueGray,colframe=RuleGray,boxrule=.55pt,arc=1.5mm,title={Topic Summary},fonttitle=\bfseries\color{Navy},left=2mm,right=2mm,top=1mm,bottom=1mm}
\newtcolorbox{practicebox}{enhanced,breakable,colback=white,colframe=Green,boxrule=.75pt,arc=2mm,title={Practitioner Lens},fonttitle=\bfseries}
\newtcolorbox{scopenotebox}{enhanced,breakable,colback=Warm,colframe=Risk,boxrule=.7pt,arc=2mm,title={Terminology and Scope Note},fonttitle=\bfseries}
\newtcolorbox{synthesisbox}{enhanced,breakable,colback=Navy!5,colframe=Navy,boxrule=.85pt,arc=2mm,title={Chapter Synthesis},fonttitle=\bfseries}
\newtcolorbox{takeawaybox}{enhanced,breakable,colback=Teal!8,colframe=Teal,boxrule=.9pt,arc=2mm,title={One-Sentence Takeaway},fonttitle=\bfseries}
\newcommand{\CornellHeader}{\rowcolor{Navy}\color{white}\bfseries Cue / Active Recall & \color{white}\bfseries Notes / Reasoning \\\endhead}
\newcommand{\CornellRow}[2]{\cellcolor{CueGray}\textbf{#1} & #2 \\\addlinespace[3pt]}

\newcommand{\ChapterMark}{Chapter 2}
\title{Chapter 2: Case Study Background\\[2mm]\large Cornell Notes}
\author{}
\date{}
\hypersetup{pdftitle={Chapter 2: Case Study Background - Cornell Notes},pdfauthor={},pdfsubject={Cloud computing and architecture study notes},pdfkeywords={cloud, architecture, security, operations, Cornell notes}}
\begin{document}
\maketitle
\thispagestyle{fancy}
\vspace{-1.5em}
\begin{orientationbox}
\textbf{Central problem.} Compare three organizations whose different starting conditions make cloud adoption, provider operation, and provider migration distinct architecture problems.\par\smallskip
\textbf{Prerequisites.} Basic cloud terminology, application portfolios, data centers, service delivery, risk, cost, and project planning.\par\smallskip
\textbf{Connection.} The neutral cases recur as practical tests for mechanisms, architectures, security choices, pricing decisions, and service-level commitments.
\end{orientationbox}
\section{Learning Objectives}
\begin{enumerate}
\item Define the principal concepts and decision boundaries in Case Study Background.
\item Explain how 2.1 Case A - enterprise portfolio transformation contributes to the chapter's central problem.
\item Distinguish Existing environment from Business objective using responsibility, behavior, and evidence.
\item Trace components, data, control, identity, capacity, and failure across 2.2 Case B - infrastructure operator becoming a provider.
\item Analyze benefits together with assumptions, limitations, security risks, and operational tradeoffs.
\item Evaluate provider and consumer responsibilities at every relevant trust or administrative boundary.
\item Select verification evidence for architecture, configuration, monitoring, cost, and service-quality claims.
\item Apply the chapter's reasoning to a realistic design, review, migration, or incident scenario.
\end{enumerate}
\section{Cornell Cue-and-Notes}
\Needspace{9\baselineskip}
\subsection{2.1 Case A - enterprise portfolio transformation}
\begin{longtable}{@{}C N@{}}
\CornellHeader
\CornellRow{What should be remembered about Existing environment?}{Multiple acquisitions left similar applications running on different platforms with little portfolio consolidation. Outsourcing reduced some labor expense but did not remove duplicated functionality or architectural complexity. \textbf{Review lens.} Inventory applications by capability, criticality, usage, users, technology, dependencies, data sensitivity, and retirement constraints before selecting migration candidates.}
\CornellRow{What should be remembered about Business objective?}{The organization seeks sustainable cost reduction, operational efficiency, agility, scalability, and reliability. These outcomes require redesign of the application landscape rather than a simple hosting-location change. \textbf{Review lens.} Define baselines and target measures for cost, delivery speed, capacity, reliability, and operational effort so benefits can be tested after migration.}
\CornellRow{What should be remembered about Roadmap?}{The staged approach screens the portfolio, performs deeper candidate assessment, designs a target architecture, builds a business case, and pilots a low-risk workload on a managed application platform. \textbf{Review lens.} Use the pilot to test transformation effort, legacy integration, deployment processes, security responsibilities, cost assumptions, and rollback options.}
\CornellRow{What should be remembered about Primary risks?}{The organization is concerned about provider selection, immature offerings, loss of governance control, policy and regulatory compliance, legacy integration, and migration economics. \textbf{Review lens.} Assign each risk an owner, acceptance criterion, evidence source, treatment, contingency, and decision date.}
\end{longtable}
\begin{sectionsummarybox}
Case A is a large telecommunications supplier whose acquisitions produced overlapping applications, heterogeneous technology, rising maintenance cost, and difficult consolidation.
\end{sectionsummarybox}
\Needspace{9\baselineskip}
\subsection{2.2 Case B - infrastructure operator becoming a provider}
\begin{longtable}{@{}C N@{}}
\CornellHeader
\CornellRow{What should be remembered about Existing environment?}{The operator has substantial compute, storage, network redundancy, and data-center expertise. Hardware standardization has improved, but software diversity and customer-specific configurations keep operations and contracts complex. \textbf{Review lens.} Measure configuration variance, manual fulfillment effort, capacity utilization, service dependencies, contract variants, and the portion of infrastructure eligible for pooling.}
\CornellRow{What should be remembered about Business objective?}{The strategy aims to standardize offerings, improve cost effectiveness, optimize operations, diversify the service portfolio, and expose repeatable self-service infrastructure capabilities. \textbf{Review lens.} Define a service catalog with explicit eligibility, tenancy, support, security, metering, pricing, and service-quality rules.}
\CornellRow{What should be remembered about Roadmap?}{The operator selects an infrastructure service model, introduces a uniform platform across two data centers, identifies pilot consumers, establishes initial contracts and service levels, and evaluates operation before broad release. \textbf{Review lens.} Gate expansion on pilot evidence for provisioning, isolation, monitoring, billing, incident response, availability, support, and consumer adoption.}
\CornellRow{What should be remembered about Primary risks?}{Long public-sector adoption cycles, legacy contract rules, regulatory uncertainty, software heterogeneity, up-front investment, and unclear pricing or service guarantees can undermine the program. \textbf{Review lens.} Treat legal and procurement readiness as architecture dependencies rather than late administrative tasks.}
\end{longtable}
\begin{sectionsummarybox}
Case B operates large, interconnected data centers and supplies customized public-sector infrastructure services; it wants to standardize its portfolio and introduce cloud-style delivery.
\end{sectionsummarybox}
\Needspace{9\baselineskip}
\subsection{2.3 Case C - established cloud service migration}
\begin{longtable}{@{}C N@{}}
\CornellHeader
\CornellRow{What should be remembered about Existing environment?}{The company depends heavily on cloud-hosted customer applications and SaaS productivity services, while retaining a relatively small on-premises estate. Its applications support rich interaction, multiple languages, social features, and sensitive use by children and parents. \textbf{Review lens.} Map user data, geographic access, identity, parental controls, latency-sensitive functions, platform dependencies, and privacy obligations before migration.}
\CornellRow{What should be remembered about Business objective?}{The organization needs more predictable scaling, stronger service levels, better economics, broader global support, and a platform that can sustain its application frameworks. \textbf{Review lens.} Translate each objective into measurable workload, latency, availability, portability, security, and cost requirements.}
\CornellRow{What should be remembered about Roadmap?}{The migration plan calls for technical and economic impact analysis, a provider-selection decision framework, portability assessment, and evaluation of current-provider cooperation. \textbf{Review lens.} Include data export, dependency replacement, parallel operation, validation, cutover, rollback, contract termination, and post-migration evidence retention.}
\CornellRow{What should be remembered about Primary risks?}{Provider-specific services may create lock-in; unpredictable demand can expose capacity limits; migration can interrupt customer service or alter security and privacy controls. \textbf{Review lens.} Prototype the most proprietary and latency-sensitive paths first, and require an exit plan before committing to the target provider.}
\end{longtable}
\begin{sectionsummarybox}
Case C already runs interactive web and mobile services in a managed application environment but has reached limits in scalability, availability, cost, and framework support.
\end{sectionsummarybox}
\Needspace{9\baselineskip}
\subsection{Cross-case comparison}
\begin{longtable}{@{}C N@{}}
\CornellHeader
\CornellRow{What should be remembered about Different transformation types?}{Case A is primarily portfolio and application transformation; Case B is operating-model and service-product transformation; Case C is portability and provider-transition transformation. \textbf{Review lens.} Choose assessment criteria and success measures that match the transformation type.}
\CornellRow{What should be remembered about Shared decision pattern?}{All three cases require current-state evidence, target outcomes, model selection, security analysis, economic justification, staged validation, and defined operating ownership. \textbf{Review lens.} Reject roadmaps that skip baseline measurement or treat a successful prototype as proof of production readiness.}
\end{longtable}
\begin{sectionsummarybox}
Cloud adoption is not one journey: Case A modernizes a fragmented estate, Case B productizes infrastructure, and Case C changes providers while preserving a live service.
\end{sectionsummarybox}
\begin{scopenotebox}
Distinctive organization names are replaced with neutral case labels while technical facts and decision patterns are preserved.
\end{scopenotebox}
\begin{practicebox}
For each model, mechanism, architecture, metric, or case decision, record the objective, owner, trust boundary, inputs, dependencies, expected behavior, failure consequence, measurement point, response, and retained evidence. A diagram or product label is a hypothesis until runtime and recovery tests support it.
\end{practicebox}
\section{Chapter Recap Matrix}
\begingroup\small\setlength{\tabcolsep}{4pt}
\begin{longtable}{@{}>{\RaggedRight\arraybackslash}p{0.235\textwidth}>{\RaggedRight\arraybackslash}p{0.37\textwidth}>{\RaggedRight\arraybackslash}p{0.325\textwidth}@{}}
\toprule\rowcolor{Navy}\color{white}\textbf{Element} & \color{white}\textbf{Purpose or Behavior} & \color{white}\textbf{Decision and Evidence}\\\midrule\endfirsthead
\toprule\rowcolor{Navy}\color{white}\textbf{Element} & \color{white}\textbf{Purpose or Behavior} & \color{white}\textbf{Decision and Evidence}\\\midrule\endhead
Existing environment & Multiple acquisitions left similar applications running on different platforms with little portfolio consolidation. & Inventory applications by capability, criticality, usage, users, technology, dependencies, data sensitivity, and retirement constraints before selecting migration candidates. \\\midrule
Business objective & The organization seeks sustainable cost reduction, operational efficiency, agility, scalability, and reliability. & Define baselines and target measures for cost, delivery speed, capacity, reliability, and operational effort so benefits can be tested after migration. \\\midrule
Roadmap & The staged approach screens the portfolio, performs deeper candidate assessment, designs a target architecture, builds a business case, and pilots a low-risk workload on a managed application platform. & Use the pilot to test transformation effort, legacy integration, deployment processes, security responsibilities, cost assumptions, and rollback options. \\\midrule
Primary risks & The organization is concerned about provider selection, immature offerings, loss of governance control, policy and regulatory compliance, legacy integration, and migration economics. & Assign each risk an owner, acceptance criterion, evidence source, treatment, contingency, and decision date. \\\midrule
Existing environment & The operator has substantial compute, storage, network redundancy, and data-center expertise. & Measure configuration variance, manual fulfillment effort, capacity utilization, service dependencies, contract variants, and the portion of infrastructure eligible for pooling. \\\midrule
Business objective & The strategy aims to standardize offerings, improve cost effectiveness, optimize operations, diversify the service portfolio, and expose repeatable self-service infrastructure capabilities. & Define a service catalog with explicit eligibility, tenancy, support, security, metering, pricing, and service-quality rules. \\\midrule
Roadmap & The operator selects an infrastructure service model, introduces a uniform platform across two data centers, identifies pilot consumers, establishes initial contracts and service levels, and evaluates operation before broad release. & Gate expansion on pilot evidence for provisioning, isolation, monitoring, billing, incident response, availability, support, and consumer adoption. \\\midrule
Primary risks & Long public-sector adoption cycles, legacy contract rules, regulatory uncertainty, software heterogeneity, up-front investment, and unclear pricing or service guarantees can undermine the program. & Treat legal and procurement readiness as architecture dependencies rather than late administrative tasks. \\\midrule
Existing environment & The company depends heavily on cloud-hosted customer applications and SaaS productivity services, while retaining a relatively small on-premises estate. & Map user data, geographic access, identity, parental controls, latency-sensitive functions, platform dependencies, and privacy obligations before migration. \\\midrule
Business objective & The organization needs more predictable scaling, stronger service levels, better economics, broader global support, and a platform that can sustain its application frameworks. & Translate each objective into measurable workload, latency, availability, portability, security, and cost requirements. \\\midrule
Roadmap & The migration plan calls for technical and economic impact analysis, a provider-selection decision framework, portability assessment, and evaluation of current-provider cooperation. & Include data export, dependency replacement, parallel operation, validation, cutover, rollback, contract termination, and post-migration evidence retention. \\\midrule
Primary risks & Provider-specific services may create lock-in; unpredictable demand can expose capacity limits; migration can interrupt customer service or alter security and privacy controls. & Prototype the most proprietary and latency-sensitive paths first, and require an exit plan before committing to the target provider. \\\midrule
Different transformation types & Case A is primarily portfolio and application transformation; Case B is operating-model and service-product transformation; Case C is portability and provider-transition transformation. & Choose assessment criteria and success measures that match the transformation type. \\\midrule
Shared decision pattern & All three cases require current-state evidence, target outcomes, model selection, security analysis, economic justification, staged validation, and defined operating ownership. & Reject roadmaps that skip baseline measurement or treat a successful prototype as proof of production readiness. \\\midrule
\bottomrule
\end{longtable}
\endgroup
\section{Practical Review Checklist}
\begin{longtable}{@{}>{\bfseries\color{Teal}}p{0.16\textwidth}>{\RaggedRight\arraybackslash}p{0.77\textwidth}@{}}
\toprule\textbf{Area} & \textbf{Verification action}\\\midrule\endhead
Scope & $\square$ Define the in-scope services, workloads, users, data, environments, regions, and dependencies for Case Study Background. \\\addlinespace[2pt]
Roles & $\square$ Name the provider, consumer, owner, administrator, broker, support, security, finance, and risk responsibilities that apply. \\\addlinespace[2pt]
Boundaries & $\square$ Map organizational, trust, network, tenancy, management, data-location, and failure-domain boundaries. \\\addlinespace[2pt]
Architecture & $\square$ Document how Existing environment, Business objective, and supporting mechanisms interact through normal, peak, maintenance, and failed states. \\\addlinespace[2pt]
Identity & $\square$ Verify human and workload identities, privileges, credentials, delegation, administrative access, revocation, and audit evidence. \\\addlinespace[2pt]
Data & $\square$ Trace data collection, transfer, processing, storage, replication, backup, retention, export, and deletion. \\\addlinespace[2pt]
Quality & $\square$ Set measurable capacity, latency, availability, reliability, scalability, and recovery criteria relevant to Roadmap. \\\addlinespace[2pt]
Security & $\square$ Challenge the trust and failure assumptions surrounding Primary risks with abuse, bypass, isolation, and compromise scenarios. \\\addlinespace[2pt]
Operations & $\square$ Confirm monitoring, alert ownership, change control, patching, drift management, incident response, and tested recovery. \\\addlinespace[2pt]
Economics & $\square$ Tie resource behavior to meters, units, rates, budgets, anomalies, full-lifecycle cost, and accountable owners. \\\addlinespace[2pt]
Contracts & $\square$ Reconcile architecture behavior with service scope, formulas, exclusions, evidence, notification, remedies, audit rights, and exit terms. \\\addlinespace[2pt]
Verification & $\square$ Preserve reproducible evidence that Existing environment and the overall design satisfy their intended outcomes under realistic load and failure. \\\addlinespace[2pt]
\bottomrule
\end{longtable}
\begin{synthesisbox}
The three cases demonstrate that cloud strategy must begin with context. A heterogeneous enterprise needs portfolio rationalization and low-risk pilots; an infrastructure operator needs service standardization and measurable contracts; an established cloud consumer needs portability, resiliency, and provider-exit planning. The migration path follows from constraints, not fashion.
\par\smallskip
\textbf{Cross-chapter application.} The neutral cases recur as practical tests for mechanisms, architectures, security choices, pricing decisions, and service-level commitments.
\end{synthesisbox}
\section{Self-Test Questions}
\begin{enumerate}
\item What central problem does Case Study Background address?
\item Define Existing environment and explain why it matters.
\item Compare Business objective with Roadmap. Which assumptions and responsibilities differ?
\item Trace the normal behavior or decision path associated with Primary risks.
\item Which evidence would you request when reviewing Existing environment?
\item How does this chapter connect to the broader cloud architecture and operating model?
\item Scenario: A team claims that Business objective is implemented but cannot show current configuration, monitoring, or failure-test evidence. What should the reviewer conclude and do next?
\item Scenario: A workload using Roadmap crosses a provider-consumer boundary and then fails during peak demand. Which responsibilities and failure conditions should be investigated?
\item Why is a named mechanism, model, or architecture insufficient evidence that its intended quality or security outcome is achieved?
\item How would you verify that a change involving Business objective improved the intended outcome without shifting risk or cost elsewhere?
\end{enumerate}
\clearpage
\section{Answer Key}
\begin{enumerate}
\item Compare three organizations whose different starting conditions make cloud adoption, provider operation, and provider migration distinct architecture problems.
\item Multiple acquisitions left similar applications running on different platforms with little portfolio consolidation. Outsourcing reduced some labor expense but did not remove duplicated functionality or architectural complexity. The review implication is: Inventory applications by capability, criticality, usage, users, technology, dependencies, data sensitivity, and retirement constraints before selecting migration candidates.
\item The organization seeks sustainable cost reduction, operational efficiency, agility, scalability, and reliability. These outcomes require redesign of the application landscape rather than a simple hosting-location change. In contrast, The staged approach screens the portfolio, performs deeper candidate assessment, designs a target architecture, builds a business case, and pilots a low-risk workload on a managed application platform. Compare control, ownership, state, dependencies, operating evidence, and failure impact rather than names alone.
\item The organization is concerned about provider selection, immature offerings, loss of governance control, policy and regulatory compliance, legacy integration, and migration economics. A trace should name initiating actor, input or signal, decision point, participating components, state change, output, monitoring, and failure handling.
\item Useful evidence should demonstrate the review lens: Measure configuration variance, manual fulfillment effort, capacity utilization, service dependencies, contract variants, and the portion of infrastructure eligible for pooling. Include architecture, configuration, event, test, and ownership evidence where applicable.
\item The neutral cases recur as practical tests for mechanisms, architectures, security choices, pricing decisions, and service-level commitments.
\item The implementation claim is unsubstantiated. Request the evidence needed to establish current state, then test positive, negative, overload, and failure paths. For this mechanism specifically: Define baselines and target measures for cost, delivery speed, capacity, reliability, and operational effort so benefits can be tested after migration.
\item Investigate the agreed responsibility split, boundary controls, capacity and quota, workload state, dependent services, provider telemetry, consumer telemetry, escalation, recovery, and whether the contract measures the consumer-visible failure.
\item A label expresses design intent, not runtime proof. Effectiveness depends on placement, configuration, inputs, state, control independence, capacity, monitoring, response, maintenance, and tested behavior during failure.
\item Establish a before-and-after measure, preserve comparable test conditions, verify the expected behavior, inspect adjacent dependencies and costs, and exercise rollback or recovery. Apply this specific lens: Define a service catalog with explicit eligibility, tenancy, support, security, metering, pricing, and service-quality rules.
\end{enumerate}
\section{Key-Term Recap}
\begin{longtable}{@{}>{\RaggedRight\arraybackslash}p{0.27\textwidth}>{\RaggedRight\arraybackslash}p{0.66\textwidth}@{}}
\toprule\rowcolor{Navy}\color{white}\textbf{Term} & \color{white}\textbf{Working meaning}\\\midrule\endhead
Existing environment & Multiple acquisitions left similar applications running on different platforms with little portfolio consolidation. Outsourcing reduced some labor expense but did not remove duplicated functionality or architectural complexity. \\\addlinespace[2pt]
Business objective & The organization seeks sustainable cost reduction, operational efficiency, agility, scalability, and reliability. These outcomes require redesign of the application landscape rather than a simple hosting-location change. \\\addlinespace[2pt]
Roadmap & The staged approach screens the portfolio, performs deeper candidate assessment, designs a target architecture, builds a business case, and pilots a low-risk workload on a managed application platform. \\\addlinespace[2pt]
Primary risks & The organization is concerned about provider selection, immature offerings, loss of governance control, policy and regulatory compliance, legacy integration, and migration economics. \\\addlinespace[2pt]
Existing environment & The operator has substantial compute, storage, network redundancy, and data-center expertise. Hardware standardization has improved, but software diversity and customer-specific configurations keep operations and contracts complex. \\\addlinespace[2pt]
Business objective & The strategy aims to standardize offerings, improve cost effectiveness, optimize operations, diversify the service portfolio, and expose repeatable self-service infrastructure capabilities. \\\addlinespace[2pt]
Roadmap & The operator selects an infrastructure service model, introduces a uniform platform across two data centers, identifies pilot consumers, establishes initial contracts and service levels, and evaluates operation before broad release. \\\addlinespace[2pt]
Primary risks & Long public-sector adoption cycles, legacy contract rules, regulatory uncertainty, software heterogeneity, up-front investment, and unclear pricing or service guarantees can undermine the program. \\\addlinespace[2pt]
Existing environment & The company depends heavily on cloud-hosted customer applications and SaaS productivity services, while retaining a relatively small on-premises estate. Its applications support rich interaction, multiple languages, social features, and sensitive use by children and parents. \\\addlinespace[2pt]
Business objective & The organization needs more predictable scaling, stronger service levels, better economics, broader global support, and a platform that can sustain its application frameworks. \\\addlinespace[2pt]
Roadmap & The migration plan calls for technical and economic impact analysis, a provider-selection decision framework, portability assessment, and evaluation of current-provider cooperation. \\\addlinespace[2pt]
Primary risks & Provider-specific services may create lock-in; unpredictable demand can expose capacity limits; migration can interrupt customer service or alter security and privacy controls. \\\addlinespace[2pt]
Different transformation types & Case A is primarily portfolio and application transformation; Case B is operating-model and service-product transformation; Case C is portability and provider-transition transformation. \\\addlinespace[2pt]
Shared decision pattern & All three cases require current-state evidence, target outcomes, model selection, security analysis, economic justification, staged validation, and defined operating ownership. \\\addlinespace[2pt]
\bottomrule
\end{longtable}
\begin{takeawaybox}
A cloud roadmap is credible only when it connects the organization's current estate, business objectives, technical constraints, risks, metrics, and pilot evidence.
\end{takeawaybox}
\end{document}
